\documentclass[10pt,a4paper]{amsart}
\usepackage{amsmath,amssymb,amsthm,graphicx}
% For commutative diagrams we use Xy-pic
\usepackage[all]{xy}

% We use 2cell for 2-commutative diagrams.
\xyoption{2cell}
% Theorem environments.
%
\theoremstyle{plain}
\newtheorem{theorem}[subsection]{Theorem}
\newtheorem{proposition}[subsection]{Proposition}
\newtheorem{lemma}[subsection]{Lemma}

\theoremstyle{definition}
\newtheorem{definition}[subsection]{Definition}
\newtheorem{example}[subsection]{Example}
\newtheorem{exercise}[subsection]{Exercise}
\newtheorem{situation}[subsection]{Situation}

\theoremstyle{remark}
\newtheorem{remark}[subsection]{Remark}
\newtheorem{remarks}[subsection]{Remarks}

\numberwithin{equation}{subsection}

% Macros
%
\def\lim{\mathop{\mathrm{lim}}\nolimits}
\def\colim{\mathop{\mathrm{colim}}\nolimits}
\def\Spec{\mathop{\mathrm{Spec}}}
\def\Hom{\mathop{\mathrm{Hom}}\nolimits}
\def\Ext{\mathop{\mathrm{Ext}}\nolimits}
\def\SheafHom{\mathop{\mathcal{H}\!\mathit{om}}\nolimits}
\def\SheafExt{\mathop{\mathcal{E}\!\mathit{xt}}\nolimits}
\def\Sch{\mathit{Sch}}
\def\Mor{\mathop{\mathrm{Mor}}\nolimits}
\def\Ob{\mathop{\mathrm{Ob}}\nolimits}
\def\Sh{\mathop{\mathit{Sh}}\nolimits}
\def\NL{\mathop{N\!L}\nolimits}
\def\CH{\mathop{\mathrm{CH}}\nolimits}
\def\proetale{{pro\text{-}\acute{e}tale}}
\def\etale{{\acute{e}tale}}
\def\QCoh{\mathit{QCoh}}
\def\Ker{\mathop{\mathrm{Ker}}}
\def\Im{\mathop{\mathrm{Im}}}
\def\Coker{\mathop{\mathrm{Coker}}}
\def\Coim{\mathop{\mathrm{Coim}}}

% Boxtimes
%
\DeclareMathSymbol{\boxtimes}{\mathbin}{AMSa}{"02}

%
% Macros for moduli stacks/spaces
%
\def\QCohstack{\mathcal{QC}\!\mathit{oh}}
\def\Cohstack{\mathcal{C}\!\mathit{oh}}
\def\Spacesstack{\mathcal{S}\!\mathit{paces}}
\def\Quotfunctor{\mathrm{Quot}}
\def\Hilbfunctor{\mathrm{Hilb}}
\def\Curvesstack{\mathcal{C}\!\mathit{urves}}
\def\Polarizedstack{\mathcal{P}\!\mathit{olarized}}
\def\Complexesstack{\mathcal{C}\!\mathit{omplexes}}
% \Pic is the operator that assigns to X its picard group, usage \Pic(X)
% \Picardstack_{X/B} denotes the Picard stack of X over B
% \Picardfunctor_{X/B} denotes the Picard functor of X over B
\def\Pic{\mathop{\mathrm{Pic}}\nolimits}
\def\Picardstack{\mathcal{P}\!\mathit{ic}}
\def\Picardfunctor{\mathrm{Pic}}
\def\Deformationcategory{\mathcal{D}\!\mathit{ef}}
\usepackage{cleveref}
\crefname{originalreference}{original reference}{original references}
\Crefname{originalreference}{Original reference}{Original references}
\makeatletter
\newlabel{lemma-proper-base-change-stalk}{{0DDF}{source}{}{sourcepointer-1}{}}
\newlabel{lemma-proper-base-change-stalk@cref}{{[originalreference][0][]0DDF}{[1][0][]source}{}{}{}}
\makeatother
\crefname{originalreference}{original reference}{original references}
\Crefname{originalreference}{Original reference}{Original references}
\makeatletter
\newlabel{lemma-specialization-cospecialization}{{0GJV}{source}{}{sourcepointer-2}{}}
\newlabel{lemma-specialization-cospecialization@cref}{{[originalreference][0][]0GJV}{[1][0][]source}{}{}{}}
\makeatother
\crefname{originalreference}{original reference}{original references}
\Crefname{originalreference}{Original reference}{Original references}
\makeatletter
\newlabel{lemma-torsion-direct-image}{{0DDD}{source}{}{sourcepointer-3}{}}
\newlabel{lemma-torsion-direct-image@cref}{{[originalreference][0][]0DDD}{[1][0][]source}{}{}{}}
\makeatother
\title{For a proper locally acyclic morphism with torsion coefficients the specialization and cospecialization maps are isomorphisms}
\author{The Stacks Project Authors}
\date{stacks-project:etale-cohomology.tex:local-snapshot-20261011; GNU FDL 1.2 or later}
\begin{document}
\maketitle
\noindent\textit{Editorial note.} All delivered mathematical text is the author's own, copied from the pinned Stacks Project chapter file etale-cohomology.tex; no statement or proof step was rewritten, reordered or removed.\par\medskip
\noindent\textit{Editorial note.} The counted proof is the author's complete argument: the proof is written out in full: the identification of both sides by proper base change, the reformulation of the cospecialization isomorphism as a pullback arrow, the verification of that arrow through the auxiliary torsion complex, and the diagram argument for the specialization map.\par\medskip
\noindent\textit{Editorial note.} The label lemma-sp-isom-proper-torsion-loc-ac names the counted statement, and the proof environment immediately following it in the source is delivered verbatim as its proof.\par\medskip
\noindent\textit{Editorial note.} The counted statement is carried by the author's own lemma environment in the chapter source.\par\medskip
\noindent\textit{Editorial note.} The statement and proof are copied from the pinned chapter file; the three cited inputs are kept as external pointers with their canonical Stacks tags.\par\medskip
\noindent\textit{Editorial note.} The preamble is the chapter's own native preamble excerpt already used for this source: the Xy-pic setup (native preamble lines 14-18) and the native theorem declarations with the mathematical macros (native preamble lines 154-217).\par\medskip
\noindent\textit{Allowed standard prerequisites (editorial).}\begin{itemize}
\item Proper base change in the stalk form. Same chapter, the stalk form of the proper base change theorem identifying (Rf\_*K)\_s with R Gamma(X\_s, K). Supplies the identification of the two sides at the geometric points s and t and the final isomorphism.
\item The diagram of specialization and cospecialization maps. Same chapter, the lemma constructing sp and cosp and the diagram relating them. Supplies the maps whose invertibility is asserted and the fourth-arrow argument at the end.
\item Torsion of the direct image. Same chapter, the lemma stating that the pushforward of a torsion complex along a morphism of schemes has torsion cohomology sheaves. Supplies that K' = Rh\_*(K|\_\{X\_t\}) is bounded below with torsion cohomology sheaves, so that proper base change applies to it.
\end{itemize}\medskip
\ifdefined\hypertarget\hypertarget{sourcepointer-1}{}\fi
\noindent\textit{Source reference.} Original reference 0DDF: sources/stacks-etale-cohomology/pinned-original-source.tex, source lines 18734--18734. Prerequisite pointer invoked at proof line(s) 19527; it names a standard fact of the same chapter that the counted argument uses, and every step of the counted proof built on it is retained verbatim in the delivered body.\par\medskip
\ifdefined\hypertarget\hypertarget{sourcepointer-2}{}\fi
\noindent\textit{Source reference.} Original reference 0GJV: sources/stacks-etale-cohomology/pinned-original-source.tex, source lines 19469--19469. Prerequisite pointer invoked at proof line(s) 19531; it names a standard fact of the same chapter that the counted argument uses, and every step of the counted proof built on it is retained verbatim in the delivered body.\par\medskip
\ifdefined\hypertarget\hypertarget{sourcepointer-3}{}\fi
\noindent\textit{Source reference.} Original reference 0DDD: sources/stacks-etale-cohomology/pinned-original-source.tex, source lines 13924--13924. Prerequisite pointer invoked at proof line(s) 19549; it names a standard fact of the same chapter that the counted argument uses, and every step of the counted proof built on it is retained verbatim in the delivered body.\par\medskip
\begin{lemma}
\label{lemma-sp-isom-proper-torsion-loc-ac}
Let $f : X \to S$ be a morphism of schemes. Let $K \in D(X_\etale)$.
Assume
\begin{enumerate}
\item $K$ is bounded below, i.e., $K \in D^+(X_\etale)$,
\item $f$ is locally acyclic relative to $K$,
\item $f$ is proper, and
\item $K$ has torsion cohomology sheaves.
\end{enumerate}
Then for every geometric point $\overline{s}$ of $S$ and every geometric
point $\overline{t}$ of $\Spec(\mathcal{O}^{sh}_{S, \overline{s}})$
both the specialization map
$sp : (Rf_*K)_{\overline{s}} \to (Rf_*K)_{\overline{t}}$
and the cospecialization map
$cosp : R\Gamma(X_{\overline{t}}, K) \to R\Gamma(X_{\overline{s}}, K)$
are isomorphisms.
\end{lemma}

\begin{proof}
By the proper base change theorem (in the form of
Lemma \ref{lemma-proper-base-change-stalk}) we have
$(Rf_*K)_{\overline{s}} = R\Gamma(X_{\overline{s}}, K)$
and similarly for $\overline{t}$. The ``correct'' proof would be
to show that the argument in
Lemma \ref{lemma-specialization-cospecialization}
shows that $sp$ and $cosp$ are inverse isomorphisms in this case.
Instead we will show directly that $cosp$ is an isomorphism.
From the discussion above we see that $cosp$ is an isomorphism
if and only if pullback by $i$
$$
R\Gamma(X \times_S \Spec(\mathcal{O}^{sh}_{S, \overline{s}}),
Rh_*(K|_{X_{\overline{t}}}))
\longrightarrow
R\Gamma(X_{\overline{s}}, i^{-1}Rh_*(K|_{X_{\overline{t}}}))
$$
is an isomorphism in $D^+(\textit{Ab})$. This is true by the proper
base change theorem for the proper morphism
$f' : X \times_S \Spec(\mathcal{O}^{sh}_{S, \overline{s}}) \to
\Spec(\mathcal{O}^{sh}_{S, \overline{s}})$
by the morphism $\overline{s} \to \Spec(\mathcal{O}^{sh}_{S, \overline{s}})$
and the complex $K' = Rh_*(K|_{X_{\overline{t}}})$. The complex
$K'$ is bounded below and has torsion cohomology sheaves
by Lemma \ref{lemma-torsion-direct-image}.
Since $\Spec(\mathcal{O}^{sh}_{S, \overline{s}})$ is strictly henselian
with $\overline{s}$ lying over the closed point,
we see that the source of the displayed arrow equals
$(Rf'_*K')_{\overline{s}}$ and the target equals
$R\Gamma(X_{\overline{s}}, K')$ and the displayed map is an isomorphism
by the already used
Lemma \ref{lemma-proper-base-change-stalk}.
Thus we see that three out of the four arrows in the diagram
of Lemma \ref{lemma-specialization-cospecialization} are isomorphisms
and we conclude.
\end{proof}

\end{document}
